\chapter{Слабые взаимодействия, вириальные поправки и неполные дифференциалы}

\lectureoverview{Мы начинаем с операционального смысла измерения и температуры, а затем спрашиваем, как разреженный газ отклоняется от закона идеального газа, когда короткодействующие взаимодействия слабо включаются. Статистическая сумма дает первые поправки к энергии и давлению, обусловленные взаимодействием, и определяет область их применимости. В заключительном математическом отступлении функции состояния отличаются от теплоты и работы, зависящих от пути, с помощью полных дифференциалов, замкнутых контуров и конкретной проверки на идеальном газе.}

\section{Что на самом деле измеряет прибор}

Ни один лабораторный прибор не обращается к природе напрямую и не выдает абстрактную математическую переменную. Пружинный динамометр регистрирует удлинение; термометр — высоту столбика жидкости. Первое мы называем измерением силы, а второе — измерением температуры лишь потому, что между видимым откликом и искомой величиной установлено воспроизводимое соотношение.

Эта связь имеет ограниченную область применимости. Пружина подчиняется закону Гука лишь пока сохраняет упругость, а достаточно большая сила в конце концов деформирует ее или ломает. Более общо, если считываемая величина \(R\) гладко зависит от физической переменной \(X\), то вблизи выбранной рабочей точки \begin{equation}
 \Delta R \simeq
 \left.\frac{dR}{dX}\right|_{X_0}\Delta X.
 \label{eq:sm6-local-linear-response}
\end{equation} Малые изменения показания пропорциональны малым изменениям искомой переменной даже тогда, когда полная кривая отклика нелинейна. Такая локальная линеаризация и составляет обычную логику калибровки.

Показание может также зависеть от посторонних переменных. Пружина может слабо зависеть от температуры или давления. Высота ртутного столбика в термометре может зависеть от ориентации через действие тяжести, от свойств стекла или от медленных изменений материалов. Процедура измерения не доказывает, что этих зависимостей нет; она либо удерживает их фиксированными, либо показывает, что в требуемом масштабе ими можно пренебречь.

\begin{classroomqa} \audiencequestion{Могут ли локальное гравитационное поле или направление, в котором держат термометр, изменить показание?} \lecturerresponse{Да, в принципе. Если ориентация оказывает заметное влияние, то указание держать термометр вертикально становится частью определения прибора. У поверхности Земли гравитационное поле достаточно стабильно, поэтому остающаяся вариация обычно пренебрежимо мала. \lecturetimestamp{00:04:08--00:05:00}} \end{classroomqa}

\begin{classroomqa} \audiencequestion{Могут ли старение или химические изменения медленно изменить калибровку?} \lecturerresponse{Могут. Ртуть и стекло достаточно инертны для обычных измерений, тогда как изменения на фантастически больших временны́х масштабах для эксперимента несущественны. Калибровка всегда есть утверждение о заданных приборе, диапазоне, протоколе и длительности. \lecturetimestamp{00:05:03--00:06:11}} \end{classroomqa}

Для термометра Цельсия удобную шкалу задают два воспроизводимых опорных состояния. Если \(h_f\) и \(h_b\) — это высоты столбика при точках замерзания и кипения воды соответственно, то в диапазоне, где отклик почти линейный, \begin{equation}
 T_{\mathrm C}=100\,\frac{h-h_f}{h_b-h_f}.
 \label{eq:sm6-celsius-calibration}
\end{equation} Выбор \(T_{\mathrm C}=0\) в точке замерзания условен. Это не состояние нулевой тепловой энергии, точно так же как нуль градусов по Фаренгейту в зимний день не является абсолютным нулем.

\begin{figure}[t] \centering
\includegraphics[width=0.90\textwidth]{lecture_06_figure_03.png} \caption{Гладкий, но нелинейный отклик термометра, с выделенным коротким интервалом для линейной калибровки.} \lectureframe{00:07:55} \end{figure}

\section{Равновесие и температура как шкала энергии}

Термометр измеряет температуру окружающей среды лишь после того, как он сам и близлежащая среда достигли теплового равновесия. До этого его показание лишь изменяется по направлению к равновесию. Практический критерий элементарен: подождать, пока показание перестанет изменяться. Простота этого правила скрывает физические процессы, определяющие время ожидания.

\begin{classroomqa} \audiencequestion{Ограничивается ли установление равновесия лишь пограничным слоем вокруг резервуара, и как долго следует ждать?} \lecturerresponse{Энергия должна пройти через стекло и обменяться с окружающей средой, поэтому время зависит от термометра, его теплопроводности и среды. Непосредственная окрестность тоже должна быть репрезентативной для однородного равновесного состояния. Экспериментально показание принимают, когда оно перестает дрейфовать. \lecturetimestamp{00:08:45--00:10:13}} \end{classroomqa}

Температура и плотность играют в этом времени ожидания разные роли. Очень разреженный газ может быть очень горячим: немногие присутствующие молекулы могут нести большие кинетические энергии. И все же эти молекулы редко ударяют по термометру, поэтому установление равновесия может идти медленно. Уменьшение плотности в сто или в тысячу раз при фиксированной температуре не уменьшает молекулярный энергетический масштаб; оно уменьшает скорость, с которой эта энергия передается прибору.

\begin{classroomqa} \audiencequestion{Может ли излучение в комнате иметь температуру, отличную от температуры воздуха?} \lecturerresponse{Да. Вещество и излучение имеют общую температуру только тогда, когда они находятся в тепловом равновесии друг с другом. Нейтральный воздух слабо взаимодействует со значительной частью оптического излучения в освещенной комнате, и популяция фотонов не обязана быть равновесным полем черного тела при температуре воздуха. \lecturetimestamp{00:11:26--00:12:20}} \end{classroomqa}

После установления равновесия температуру лучше всего рассматривать как энергетическую шкалу. Мы по-прежнему используем единицы, в которых \(k_B=1\), так что сама \(T\) имеет размерность энергии. Для одноатомного идеального газа \begin{equation}
 E_0=\frac{3}{2}NT.
 \label{eq:sm6-ideal-energy}
\end{equation} Тепловая энергия на молекулу определяется \(T\); соответствующая скорость все же зависит от массы молекулы.

\section{Включение слабого взаимодействия}

\begin{classroomqa} \audiencequestion{При каких условиях перестает работать приближение идеального газа, и можно ли показать это математически?} \lecturerresponse{Начните с решенного случая невзаимодействующего газа и включите малое взаимодействие. Разложите по этому изменению, вычислите первую поправку и сравните ее с идеальным членом. Приближение нарушается, когда поправка перестает быть малой. \lecturetimestamp{00:12:25--00:14:40}} \end{classroomqa}

Это стандартная пертурбативная стратегия. Если задача решается при \(\lambda=0\), вводят параметр связи \(\lambda\), пишут \begin{equation}
 F(\lambda)=F(0)+\lambda F_1+\lambda^2F_2+\cdots,
 \label{eq:sm6-perturbative-template}
\end{equation} и спрашивают, меньше ли каждый сохраняемый член предыдущего. Здесь опорной задачей служит идеальный газ. Физическая малость обеспечивается разреженным образцом, малым радиусом взаимодействия по сравнению со средним расстоянием и потенциалом, достаточно слабым, чтобы первый член по взаимодействию был полезен. Большая плотность или дальнодействующая сила могут сделать приближение неприменимым, даже если каждое отдельное взаимодействие кажется слабым.

Пусть \(x_i\in\mathbb R^3\) и \(p_i\in\mathbb R^3\) — положение и импульс молекулы \(i\). Для тождественных молекул массы \(m\) возьмем \begin{equation}
 H(X,P)=
 \sum_{i=1}^{N}\frac{p_i^2}{2m}
 +\sum_{i<j}u\!\left(|x_i-x_j|\right),
 \label{eq:sm6-interacting-hamiltonian}
\end{equation} где \(X=(x_1,\ldots,x_N)\) и \(P=(p_1,\ldots,p_N)\). Ограничение \(i<j\) исключает пары частицы с самой собой и учитывает \((i,j)\) только один раз: \begin{equation}
 \sum_{i<j}1=\binom{N}{2}=\frac{N(N-1)}{2}.
 \label{eq:sm6-pair-count}
\end{equation}

\begin{classroomqa} \audiencequestion{Почему не суммировать независимо по обеим меткам частиц?} \lecturerresponse{Пара \((3,7)\) — та же физическая пара, что и \((7,3)\), а частица не образует пару сама с собой. Ограничение суммы условием \(i<j\) учитывает оба обстоятельства. Обозначение массы сохраняется как \(m\), тогда как \(i,j\) нумеруют частицы, чтобы избежать конфликта обозначений, который может возникнуть, когда \(m\) используют в обеих ролях. \lecturetimestamp{00:17:52--00:19:12}} \end{classroomqa}

Запишем полный конфигурационный потенциал как \begin{equation}
 U(X)=\sum_{i<j}u\!\left(|x_i-x_j|\right).
 \label{eq:sm6-total-potential}
\end{equation} Предполагается, что взаимодействие трансляционно инвариантно и короткодействующее: \(u(r)\) становится пренебрежимо малым, когда \(r\) велик по сравнению с длиной молекулярного масштаба.

\section{Интегральная характеристика парного взаимодействия}

В первую поправку входит только один интеграл от слабого парного потенциала: \begin{equation}
 U_0\equiv\int_{\mathbb R^3}d^3r\,u(r).
 \label{eq:sm6-u0-definition}
\end{equation} Его размерность — энергия, умноженная на объем. Он объединяет типичный масштаб взаимодействия с объемом, в пределах которого оно заметно.

\begin{figure}[t] \centering
\includegraphics[width=0.90\textwidth]{lecture_06_figure_04.png} \caption{Короткодействующий парный потенциал, сведенный к интегральной характеристике \(U_0=\int d^3r\,u(r)\).} \lectureframe{00:25:15} \end{figure}

Чтобы понять, почему \(U_0\) полезен, рассмотрим одну пару: \begin{equation}
 I_2=\int_V d^3x_1\,d^3x_2\,
 u\!\left(|x_1-x_2|\right).
 \label{eq:sm6-two-particle-integral}
\end{equation} Введем относительную координату \(r=x_2-x_1\). Вдали от границы интеграл по относительной координате дает \(U_0\), тогда как положение пары как целого пробегает весь сосуд. Поэтому \begin{equation}
 I_2=V U_0+\text{surface corrections}.
 \label{eq:sm6-pair-integral-vu0}
\end{equation} Поправка возникает из-за частиц, находящихся от стенки на расстоянии не больше радиуса взаимодействия, где часть окрестности в пространстве относительной координаты лежит вне ящика. Ее относительная важность масштабируется как отношение площади поверхности к объему и исчезает в макроскопическом пределе.

\begin{classroomqa} \audiencequestion{Относится ли \(U_0\) к одной выбранной паре, и удерживается ли одна частица на фиксированном расстоянии от другой?} \lecturerresponse{Оно относится к общему закону взаимодействия для любой пары. Зафиксируйте первую частицу и проинтегрируйте вторую по всем относительным положениям. Расстояние при этом интегрировании не фиксировано. Существенный вклад дает только близкая область взаимодействия, что и приводит к \(U_0\). \lecturetimestamp{00:26:03--00:27:33}} \end{classroomqa}

\begin{classroomqa}
\audiencequestion{Почему второй интеграл по координатам дает полный объем сосуда, а первый — нет?}
\lecturerresponse{После того как относительная конфигурация проинтегрирована в \(U_0\), всю пару можно переносить по всему образцу. Этот общий перенос и дает большой множитель \(V\). Область взаимодействия — это объем молекулярного масштаба, уже содержащийся в \(U_0\), а не еще один множитель, равный объему сосуда.
\lecturetimestamp{00:27:48--00:29:21}}
\end{classroomqa}

\begin{classroomqa}
\audiencequestion{Может ли значение зависеть от того, какая точка удерживается фиксированной, особенно вблизи стенки?}
\lecturerresponse{Трансляционная инвариантность делает объемное значение независимым от фиксированной точки. Вблизи стенки интеграл немного отличается, потому что некоторые соседние положения недоступны. Затрагивается лишь поверхностный слой молекулярной толщины, поэтому ошибкой можно пренебречь, когда сосуд макроскопичен.
\lecturetimestamp{00:29:24--00:31:22}}
\end{classroomqa}

\begin{classroomqa}
\audiencequestion{Является ли результат потенциальной энергией, проинтегрированной по всем двухчастичным конфигурациям, и почему он пропорционален только одному большому объему?}
\lecturerresponse{Да. Его размерность — энергия, умноженная на квадрат объема, поскольку \(V U_0\) содержит один объем сосуда и один молекулярный объем взаимодействия. Лишь первый из них растет с размером образца.
\lecturetimestamp{00:31:31--00:32:04}}
\end{classroomqa}

\begin{classroomqa}
\audiencequestion{Изменит ли выбор сферы, которая при фиксированном объеме минимизирует площадь, объемный результат?}
\lecturerresponse{Он изменяет только малую граничную поправку. Для любой обычной макроскопической формы число частиц, находящихся на молекулярном расстоянии от поверхности, пренебрежимо мало по сравнению с числом частиц в объеме. Поверхность становится важной лишь тогда, когда сам сосуд приближается по размеру к молекулярным масштабам.
\lecturetimestamp{00:32:04--00:33:58}}
\end{classroomqa}

\begin{classroomqa}
\audiencequestion{Не упущены ли столкновения при замене траекторий частиц потенциалом?}
\lecturerresponse{Нет. Сила есть производная потенциала, и именно она порождает столкновения. Равновесная статистическая механика не требует, чтобы мы отслеживали эти столкновения по одному; их эффект закодирован в энергии и, следовательно, в больцмановском весе.
\lecturetimestamp{00:33:59--00:35:29}}
\end{classroomqa}

\begin{classroomqa}
\audiencequestion{Что теряется, если сохранить парные силы, но не учитывать подлинные трехчастичные взаимодействия?}
\lecturerresponse{Микроскопическая энергия может содержать подлинные трехчастичные члены, а более высокие порядки парного потенциала также порождают многочастичные корреляции. В разреженном газе, однако, трем заданным частицам значительно менее вероятно, чем двум, оказаться в одной молекулярной окрестности. Поэтому парные эффекты дают первую поправку по плотности; кластеры более высокого порядка входят при более высоких степенях плотности.
\lecturetimestamp{00:35:29--00:38:14}}
\end{classroomqa}

\section{Статистическая сумма на автопилоте}

Преимущество статистического метода теперь видно. Как только гамильтониан задан, нет нужды угадывать подробное распределение скоростей у стенок или во время столкновений. Мы записываем каноническую статистическую сумму и прослеживаем ее алгебру:
\begin{equation}
\begin{aligned}
 Z&=
 \frac{1}{N!h^{3N}}
 \int d^{3N}P\,d^{3N}X\\
 &\qquad{}\times e^{-\beta K(P)}e^{-\beta U(X)},\\
 K(P)&=\sum_i\frac{p_i^2}{2m}.
\end{aligned}
 \label{eq:sm6-full-partition-function}
\end{equation}
Множитель \(h^{3N}\) делает классическую меру в фазовом пространстве безразмерной, а \(N!\) устраняет переучет перестановок тождественных частиц.

\begin{figure}[t]
\centering
\includegraphics[width=0.90\textwidth]{lecture_06_figure_05.png}
\caption{Энергия многочастичной системы и соответствующие кинетический и конфигурационный множители в статистической сумме.}
\lectureframe{00:40:35}
\end{figure}

Импульсы и координаты разделяются. Умножение и деление конфигурационного интеграла на \(V^N\) позволяет выделить результат для идеального газа:
\begin{align}
 Z
 &=\left[
 \frac{V^N}{N!h^{3N}}
 \int d^{3N}P\,e^{-\beta K(P)}
 \right] \notag\\
 &\quad\times\left[
 \frac{1}{V^N}
 \int d^{3N}X\,e^{-\beta U(X)}
 \right] \notag\\
 &\equiv Z_0\,C(\beta,V,N).
 \label{eq:sm6-z-factorization}
\end{align}
Знакомый гауссовский множитель равен
\begin{equation}
 \int d^{3N}P\,e^{-\beta K(P)}
 =\left(\frac{2\pi m}{\beta}\right)^{3N/2}.
 \label{eq:sm6-gaussian-factor}
\end{equation}

\begin{figure}[t]
\centering
\includegraphics[width=0.90\textwidth]{lecture_06_figure_06.png}
\caption{Статистическая сумма идеального газа, выделенная из нормированного конфигурационного множителя.}
\lectureframe{00:43:45}
\end{figure}

Множитель \(N!\) не влияет ни на нормированное отношение \(C\), ни на используемые ниже конкретные производные при фиксированном \(N\), но в термодинамике в целом от него нельзя отказаться: он необходим для правильного подсчета состояний и экстенсивной энтропии.

В первом порядке по взаимодействию,
\begin{equation}
 e^{-\beta U(X)}
 =1-\beta U(X)+O(U^2).
 \label{eq:sm6-boltzmann-expansion}
\end{equation}
Следовательно,
\begin{align}
 C
 &=
 1-\frac{\beta}{V^N}
 \sum_{i<j}
 \int d^{3N}X\,u(|x_i-x_j|)
 +O(U^2) \notag\\
 &=
 1-\frac{\beta}{V^N}
 \frac{N(N-1)}{2}
 V^{N-2}(V U_0)+O(U^2) \notag\\
 &=
 1-\beta\frac{N(N-1)}{2V}U_0+O(U^2).
 \label{eq:sm6-configurational-correction}
\end{align}
У каждой пары один и тот же \emph{интеграл}, а не одно и то же мгновенное расстояние или потенциальная энергия. Каждый из оставшихся \(N-2\) интегралов по координатам дает по одному множителю \(V\).

\begin{figure}[t]
\centering
\includegraphics[width=0.90\textwidth]{lecture_06_figure_07.png}
\caption{Подсчет пар, объемы по координатам остальных частиц и первая конфигурационная поправка \(1-\beta N^2U_0/(2V)\).}
\lectureframe{00:53:40}
\end{figure}

\begin{classroomqa}
\audiencequestion{Предполагает ли главный член \(1\), что потенциал обращается в нуль при нулевом расстоянии?}
\lecturerresponse{Нет. Это член нулевого порядка в разложении экспоненты, не зависящий от значения \(u(0)\). Потенциал входит в следующий член, \(-\beta U(X)\).
\lecturetimestamp{00:46:19--00:46:48}}
\end{classroomqa}

\begin{classroomqa}
\audiencequestion{Почему каждую пару можно заменить одним и тем же выражением, если у разных пар разные расстояния?}
\lecturerresponse{Их поточечные значения друг к другу не приравниваются. Каждый член интегрируется по одному и тому же полному конфигурационному пространству, а метки частиц не отличают одну пару от другой. Интегралы имеют одинаковую форму и одинаковое значение.
\lecturetimestamp{00:50:00--00:51:07}}
\end{classroomqa}

\begin{classroomqa}
\audiencequestion{Число пар равно \(N(N-1)/2\) или
\(N(N+1)/2\)?}
\lecturerresponse{Это \(N(N-1)/2\): после выбора одной частицы остается лишь \(N-1\) различных партнеров, а деление на два устраняет обратный порядок. Для макроскопических \(N\) замена \(N(N-1)\) на \(N^2\) меняет только относительную поправку порядка \(1/N\).
\lecturetimestamp{00:48:33--00:49:56}}
\end{classroomqa}

Поскольку здесь \(\log(1-x)=-x-x^2/2-x^3/3-\cdots\), а не знакопеременный ряд, логарифм в первом порядке равен
\begin{equation}
 \log Z
 =
 \log Z_0
 -\beta\frac{N(N-1)}{2V}U_0
 +O(U^2).
 \label{eq:sm6-logz-correction}
\end{equation}
При больших \(N\) можно заменить \(N(N-1)\) на \(N^2\) там, где различие при конечном \(N\) учитывать не нужно.

\begin{remark}
Контролируемая формулировка при низкой плотности использует функцию Майера \(f(r)=e^{-\beta u(r)}-1\) и второй вириальный коэффициент
\begin{equation}
 B_2(T)=-\frac{1}{2}\int d^3r\,f(r).
 \label{eq:sm6-second-virial}
\end{equation}
Для слабого потенциала \(f(r)\simeq-\beta u(r)\), так что
\(B_2\simeq\beta U_0/2\), и тем самым получается приведенный ниже результат. Эта форма нужна для жесткого отталкивательного ядра, где прямое разложение \(e^{-\beta u}\) недопустимо.
\end{remark}

\section{Энергия, давление и смысл поправки}

Логарифм \(Z\) полезен, потому что его производные можно измерить. Предполагая, что микроскопический потенциал не зависит от температуры,
\begin{align}
 E
 &=-\left(\frac{\partial\log Z}{\partial\beta}\right)_{V,N} \notag\\
 &=
 \frac{3}{2}NT
 +\frac{N(N-1)}{2V}U_0+O(U^2) \notag\\
 &\simeq
 N\left(\frac{3}{2}T+\frac{1}{2}\rho U_0\right),
 \qquad
 \rho=\frac{N}{V}.
 \label{eq:sm6-energy-correction}
\end{align}
Кажущаяся неэкстенсивность \(N^2\) сопровождается множителем \(1/V\). При фиксированной плотности \(N^2/V=N\rho\), так что полная энергия остается пропорциональной размеру системы. Дополнительная энергия на частицу пропорциональна плотности, потому что выбранная молекула приобретает парную энергию только тогда, когда поблизости находится другая молекула.

\begin{figure}[t]
\centering
\includegraphics[width=0.90\textwidth]{lecture_06_figure_08.png}
\caption{Первая поправка по взаимодействию, переписанная как экстенсивная энергия \(N[3T/2+\rho U_0/2]\).}
\lectureframe{01:03:25}
\end{figure}

Давление измерять легче. Используя результат предыдущей главы,
\begin{equation}
 P=T\left(\frac{\partial\log Z}{\partial V}\right)_{T,N},
 \label{eq:sm6-pressure-from-z}
\end{equation}
получаем
\begin{align}
 P
 &=
 \frac{NT}{V}
 +T\frac{\partial}{\partial V}
 \left(-\beta\frac{N(N-1)}{2V}U_0\right)
 +O(U^2) \notag\\
 &=
 \rho T+\frac{1}{2}\rho^2U_0
 +O\!\left(\frac{\rho U_0}{N},U^2\right).
 \label{eq:sm6-pressure-correction}
\end{align}

\begin{figure}[t]
\centering
\includegraphics[width=0.90\textwidth]{lecture_06_figure_09.png}
\caption{Давление из производной Гельмгольца, включая первую квадратичную по плотности поправку.}
\lectureframe{01:09:50}
\end{figure}

\begin{classroomqa}
\audiencequestion{Почему именно \(\rho^2\), а не другая степень, является естественной первой поправкой?}
\lecturerresponse{Идеальный член требует присутствия одной частицы в малой области и поэтому масштабируется как \(\rho\). Для взаимодействия нужно, чтобы две частицы занимали одну молекулярную окрестность, вероятность чего масштабируется как \(\rho^2\). Множитель \(1/2\) — это тот же самый коэффициент подсчета пар, который не позволяет считать \((i,j)\) и \((j,i)\) по отдельности.
\lecturetimestamp{01:24:27--01:26:00}}
\end{classroomqa}

При фиксированной температуре можно изменять \(V\), измерять \(\rho=N/V\) и строить график давления. Идеальный закон дает линейный член. Кривизна \(P(\rho)\) обнаруживает первую поправку по взаимодействию. В рамках этого приближения слабого потенциала поправка не имеет явной зависимости от температуры, потому что множитель \(T\) в \eqref{eq:sm6-pressure-from-z} сокращает
\(\beta=1/T\).

\begin{classroomqa}
\audiencequestion{Насколько большой может стать плотность, прежде чем придется учитывать трехчастичные эффекты?}
\lecturerresponse{Подсчет по степеням плотности показывает, что следующий вириальный член пропорционален \(\rho^3\), но его коэффициент нельзя вывести только из \(U_0\). В него вносят вклад коррелированные произведения парных взаимодействий и, если они есть, подлинные трехчастичные силы. Нужно вычислить следующий кластерный коэффициент и сравнить его с членом \(\rho^2\); предполагаемую формулу не следует выдавать за результат.
\lecturetimestamp{01:10:50--01:13:38}}
\end{classroomqa}

Член идеального газа доминирует, когда величина поправки мала:
\begin{equation}
 \begin{aligned}
 \frac{1}{2}\rho^2|U_0|&\ll \rho T,\\
 |\rho U_0|&\ll T
 \qquad\text{(parametrically)}.
 \end{aligned}
 \label{eq:sm6-ideal-validity}
\end{equation}
Обе стороны второго неравенства имеют размерность энергии. Величина \(\rho U_0\) задает масштаб энергии взаимодействия на частицу, тогда как \(T\) задает масштаб кинетической энергии. Обычный разреженный газ при комнатной температуре находится далеко внутри этого режима; существенные отклонения возникают лишь тогда, когда молекулярные окрестности заселяются настолько часто, что взаимодействия начинают конкурировать с тепловым движением.

\begin{classroomqa}
\audiencequestion{Можно ли оценить этот критерий для воздуха в этой комнате?}
\lecturerresponse{Оцените \(U_0\) как энергию межмолекулярного взаимодействия, умноженную на молекулярный объём, умножьте на численную плотность и сравните результат с \(k_BT\). Грубые молекулярные оценки показывают, что в обычном комнатном воздухе масштаб взаимодействия намного меньше. Чтобы молекулы часто оказывались в пределах радиуса взаимодействия друг друга, требуется значительное сжатие.
\lecturetimestamp{01:16:24--01:18:54}}
\end{classroomqa}

\begin{classroomqa}
\audiencequestion{Может ли \(U_0\) иметь любой знак?}
\lecturerresponse{Да. В этом приближении положительный интеграл взаимодействия в сумме даёт отталкивание и повышает давление по сравнению с \(\rho T\). Отрицательный — в сумме даёт притяжение и понижает его. Эта проверка знака служит непосредственной физической проверкой \eqref{eq:sm6-pressure-correction}.
\lecturetimestamp{01:18:56--01:20:29}}
\end{classroomqa}

\begin{figure}[t]
\centering
\includegraphics[width=0.90\textwidth]{lecture_06_figure_10.png}
\caption{Качественные профили отталкивающего и притягивающего парного потенциала рядом с формулой для давления, поправка в которой имеет тот же знак, что и \(U_0\).}
\lectureframe{01:21:30}
\end{figure}

\begin{classroomqa}
\audiencequestion{Можно ли подставить в этот расчёт электростатическое отталкивание, гравитационное притяжение или реальные молекулярные силы?}
\lecturerresponse{Микроскопический потенциал задаётся атомной, молекулярной или квантовой физикой. Для нейтральных молекул обычно характерны крутое короткодействующее отталкивающее ядро и более длинный притягивающий хвост Ван-дер-Ваальса, который часто схематически представляют потенциалом 12--6. Для такого жёсткого ядра следует использовать полную формулу через функцию Майера. Гравитация иная: она дальнодействующая, её пространственный интеграл не сходится требуемым образом, и обычный короткодействующий термодинамический предел, используемый здесь, не работает.
\lecturetimestamp{01:20:40--01:23:03}}
\end{classroomqa}

Этот расчёт иллюстрирует полезное разделение труда. Микроскопическая физика задаёт \(u(r)\). Статистическая механика превращает эту энергию в \(Z\), производные от \(Z\) дают макроскопические наблюдаемые, а физические рассуждения проверяют их знак, зависимость от плотности, экстенсивность и область применимости. Формальные манипуляции становятся механическими лишь после того, как выбраны правильные энергия и приближение.

\section{Точные дифференциалы и зависимость от пути}

Теперь мы меняем математический язык. В выражении \emph{точный дифференциал} слово «точный» не означает численной точности. Оно означает, что 1-форма является дифференциалом подлинной скалярной функции.

Пусть \(f(x,y)\) обозначает высоту в точке некоторого ландшафта. Тогда при малом смещении она изменяется на
\begin{equation}
 df=
 \frac{\partial f}{\partial x}\,dx+
 \frac{\partial f}{\partial y}\,dy
 =f_x\,dx+f_y\,dy.
 \label{eq:sm6-total-differential}
\end{equation}
Предположим теперь, что нам просто дана
\begin{equation}
 \omega=F_x(x,y)\,dx+F_y(x,y)\,dy.
 \label{eq:sm6-one-form}
\end{equation}
Существует ли скаляр \(f\) такой, что \(\omega=df\)? Равенство смешанных производных требует
\begin{equation}
 \frac{\partial F_y}{\partial x}
 =
 \frac{\partial F_x}{\partial y}.
 \label{eq:sm6-exactness-condition}
\end{equation}
На односвязной области это условие нулевого ротора также является достаточным. Тогда компоненты \((F_x,F_y)\) образуют градиент \(f\).

\begin{figure}[t]
\centering
\includegraphics[width=0.90\textwidth]{lecture_06_figure_11.png}
\caption{Высота как скалярная функция, её дифференциал и условие точности в виде равенства смешанных производных.}
\lectureframe{01:34:30}
\end{figure}

То же условие имеет и интегральный смысл. Если \(\omega=df\), то
\begin{equation}
 \int_{\gamma:1\to2}\omega=f(2)-f(1)
 \label{eq:sm6-path-independent-integral}
\end{equation}
зависит только от конечных точек, а не от пути \(\gamma\). В частности,
\begin{equation}
 \oint_\gamma\omega=0
 \label{eq:sm6-closed-loop-zero}
\end{equation}
для любого стягиваемого замкнутого контура.

Высота обладает этим свойством. Если ехать между одними и теми же двумя точками по ровной дороге или через холм, суммарное изменение высоты будет одинаковым. Расход топлива таким свойством не обладает. Маршрут по холмистой местности и маршрут по ровной дороге, вообще говоря, расходуют разные количества топлива, и замкнутая поездка не возвращает сожжённое топливо. Количество оставшегося топлива, конечно, является хорошо определённой величиной в расширенном описании состояния, но оно не есть функция одного лишь положения автомобиля \((x,y)\); оно хранит память о маршруте. Усталость даёт аналогию того же рода: спуск с холма, вообще говоря, не устраняет усталости от подъёма на него.

\begin{classroomqa}
\audiencequestion{Является ли точность той же идеей, что и консервативная сила, и связана ли она с кривизной?}
\lecturerresponse{Силовое поле с нулевым ротором локально консервативно и допускает потенциал, так что терминология тесно связана. Параллельный перенос и кривизна тоже связывают изменения при обходе замкнутого контура с локальными геометрическими данными, но это более широкая тема. Здесь нужно только элементарное утверждение: нулевой ротор означает локальную независимость от пути при обычном глобальном условии на область.
\lecturetimestamp{01:39:22--01:41:07}}
\end{classroomqa}

Два примера делают критерий наглядным. Сначала возьмём
\begin{equation}
 F_x=y,\qquad F_y=x.
 \label{eq:sm6-exact-example}
\end{equation}
Обе смешанные производные равны \(1\), и действительно
\begin{equation}
 \omega=y\,dx+x\,dy=d(xy).
 \label{eq:sm6-exact-example-potential}
\end{equation}
Теперь изменим один знак:
\begin{equation}
 F_x=y,\qquad F_y=-x.
 \label{eq:sm6-inexact-example}
\end{equation}
Тогда
\begin{equation}
 \frac{\partial F_x}{\partial y}=1,
 \qquad
 \frac{\partial F_y}{\partial x}=-1,
 \label{eq:sm6-inexact-cross-derivatives}
\end{equation}
так что не существует скаляра \(f(x,y)\), градиент которого имел бы эти две компоненты.

\begin{figure}[t]
\centering
\includegraphics[width=0.90\textwidth]{lecture_06_figure_12.png}
\caption{Точный пример, порождённый \(f=xy\), и неточный пример, для которого смешанные производные равны \(1\) и \(-1\).}
\lectureframe{01:44:50}
\end{figure}

\section{Теплота и работа не являются функциями состояния}

Для простой сжимаемой системы с фиксированным числом частиц состояние определяется двумя независимыми равновесными переменными. Можно использовать \((T,V)\), \((E,V)\) или, что особенно естественно для фундаментального соотношения, \((S,V)\). Тогда
\begin{equation}
 E=E(S,V),
 \end{equation}
а соответствующий точный дифференциал имеет вид
\begin{equation}
 dE=
 \left(\frac{\partial E}{\partial S}\right)_V dS+
 \left(\frac{\partial E}{\partial V}\right)_S dV
 =T\,dS-P\,dV.
 \label{eq:sm6-fundamental-relation}
\end{equation}

\begin{classroomqa}
\audiencequestion{Означает ли «адиабатический»: медленный, теплоизолированный, обратимый или с постоянной энтропией?}
\lecturerresponse{В разных контекстах этот термин употребляется по-разному. Для рассматриваемого здесь равновесного пути изменение является квазистатическим и теплоизолированным, а потому обратимым и изоэнтропическим. Просто теплоизолированный необратимый процесс всё равно может производить энтропию, тогда как квантовая адиабатическая теорема подчёркивает медленное изменение параметров.
\lecturetimestamp{01:47:13--01:47:54}}
\end{classroomqa}

При фиксированной энтропии расширение изменяет энергию газа на \(-P\,dV\). При фиксированном объёме обратимый подвод теплоты изменяет её на \(T\,dS\). Поэтому полезно записать
\begin{equation}
 \begin{aligned}
 \delta W_{\mathrm{on}}&=-P\,dV,
 &\delta Q_{\mathrm{rev}}&=T\,dS,\\
 dE&=\delta Q_{\mathrm{rev}}+\delta W_{\mathrm{on}}.&&
 \end{aligned}
 \label{eq:sm6-first-law-increments}
\end{equation}
Символ \(\delta\), а не \(d\), заранее указывает, что теплота и работа — это приращения, зависящие от пути, а не дифференциалы функций состояния.

\begin{figure}[t]
\centering
\includegraphics[width=0.90\textwidth]{lecture_06_figure_13.png}
\caption{Первый закон для поршневой системы:
\(dE=-P\,dV+T\,dS\).}
\lectureframe{01:50:28}
\end{figure}

\begin{classroomqa}
\audiencequestion{Представляет ли \(-P\,dV\) работу, совершаемую над газом, или работу, совершаемую газом?}
\lecturerresponse{При соглашении о знаках в
\eqref{eq:sm6-first-law-increments}, это работа, совершаемая над газом. При расширении \(dV>0\) газ совершает положительную работу над поршнем, поэтому работа, совершаемая над газом, отрицательна, и газ теряет энергию. Эти два описания равны по величине и противоположны по знаку.
\lecturetimestamp{01:51:23--01:52:27}}
\end{classroomqa}

Равенство \(\delta Q_{\mathrm{rev}}=T\,dS\) нельзя в неизменном виде распространять на произвольный необратимый нагрев. Вообще говоря,
\begin{equation}
 dS=\frac{\delta Q}{T_{\mathrm{boundary}}}+dS_{\mathrm{prod}},
 \qquad dS_{\mathrm{prod}}\ge 0.
 \label{eq:sm6-entropy-production}
\end{equation}
Приведённый ниже аргумент о точном дифференциале относится к квазистатическим равновесным путям, для которых уместно именно обратимое соотношение.

\begin{classroomqa}
\audiencequestion{Предполагается ли в рассуждении о пространстве состояний, что все пути проходят через равновесные состояния?}
\lecturerresponse{Да. Операции выполняются достаточно медленно, чтобы вдоль пути оставались хорошо определёнными \(P\), \(T\), \(E\), \(S\) и \(V\). Зависимость теплоты и работы от пути не требует, чтобы систему уводили далеко от равновесия.
\lecturetimestamp{01:54:30--01:54:39}}
\end{classroomqa}

Чтобы проверить условие точности для теплоты, используем \((E,V)\) в качестве координат и перепишем первый закон:
\begin{equation}
 \delta Q_{\mathrm{rev}}=dE+P(E,V)\,dV.
 \label{eq:sm6-heat-one-form}
\end{equation}
Если бы существовала функция состояния \(Q(E,V)\) такая, что \(dQ=\delta Q_{\mathrm{rev}}\),
её частные производные должны были бы иметь вид
\begin{equation}
 \left(\frac{\partial Q}{\partial E}\right)_V=1,
 \qquad
 \left(\frac{\partial Q}{\partial V}\right)_E=P.
 \label{eq:sm6-hypothetical-heat-derivatives}
\end{equation}
Тогда равенство смешанных производных требовало бы
\begin{equation}
 \begin{aligned}
 \frac{\partial}{\partial V}
 \left(\frac{\partial Q}{\partial E}\right)_V
 &=\frac{\partial}{\partial E}
 \left(\frac{\partial Q}{\partial V}\right)_E,\\
 0&=\left(\frac{\partial P}{\partial E}\right)_V.
 \end{aligned}
 \label{eq:sm6-heat-exactness-requirement}
\end{equation}

\begin{classroomqa}
\audiencequestion{Это в точности тот же критерий ротора, что и для
\(F_x\,dx+F_y\,dy\)?}
\lecturerresponse{Да. Коэффициенты при \(dE\) и \(dV\) равны \(1\) и
\(P(E,V)\). Их смешанные производные равны \(0\) и
\((\partial P/\partial E)_V\). Теплота была бы точным дифференциалом только в том случае, если бы эти производные совпадали.
\lecturetimestamp{01:54:39--01:54:52}}
\end{classroomqa}

\begin{figure}[t]
\centering
\includegraphics[width=0.90\textwidth]{lecture_06_figure_14.png}
\caption{1-форма теплоты \(dE+P\,dV\) и те смешанные производные, которые потребовались бы, если бы она была дифференциалом функции состояния.}
\lectureframe{01:58:10}
\end{figure}

Это требование уже физически неправдоподобно: добавление энергии газу при фиксированном объёме, вообще говоря, меняет его давление. Для идеального газа можно выполнить прямой расчёт:
\begin{equation}
 PV=NT,
 \qquad
 E=\frac{3}{2}NT
 \quad\Longrightarrow\quad
 P=\frac{2E}{3V}.
 \label{eq:sm6-ideal-pressure-energy}
\end{equation}
Следовательно
\begin{equation}
 \left(\frac{\partial P}{\partial E}\right)_V
 =\frac{2}{3V}\ne0.
 \label{eq:sm6-heat-curl-fails}
\end{equation}
Критерий ротора не выполняется: \(\delta Q_{\mathrm{rev}}\) не является точным дифференциалом. Приращение работы \(\delta W_{\mathrm{on}}=-P\,dV\) также не является точным дифференциалом по той же причине.

\begin{figure}[t]
\centering
\includegraphics[width=0.90\textwidth]{lecture_06_figure_15.png}
\caption{Соотношения для идеального газа, объединённые в
\(P=2E/(3V)\).}
\lectureframe{02:00:00}
\end{figure}

\clearpage
\begin{figure}[H]
\centering
\includegraphics[width=0.90\textwidth]{lecture_06_figure_16.png}
\caption{Завершённый аргумент: поскольку
\((\partial P/\partial E)_V\ne0\), 1-форма теплоты не может быть точной.}
\lectureframe{02:00:35}
\end{figure}

Для замкнутого равновесного цикла,
\begin{equation}
 \oint dE=0,
 \qquad
 \oint\delta Q_{\mathrm{rev}}
 =-\oint\delta W_{\mathrm{on}}.
 \label{eq:sm6-cycle-balance}
\end{equation}
Система возвращается к той же энергии, но теплота и работа не обязаны по отдельности обращаться в нуль. Их значения зависят от маршрута в пространстве состояний. Вот почему циклическая машина может поглощать теплоту и совершать суммарную работу, хотя её рабочее вещество возвращается в исходное термодинамическое состояние.

\begin{classroomqa}
\audiencequestion{Связана ли эта зависимость от пути с КПД тепловой машины, гистерезисом и памятью?}
\lecturerresponse{В циклах тепловых машин она используется непосредственно: в замкнутом цикле подведённое тепло уравновешивает чистую совершённую работу, тогда как функция состояния \(E\) возвращается к своему исходному значению. Гистерезис и память также связаны с зависимостью от пути, хотя для них часто требуются дополнительные внутренние переменные или неравновесная физика. Данный вывод более узок и уже справедлив для обратимых равновесных путей: теплота и работа — это величины переноса, а не свойства, хранящиеся в термодинамическом состоянии.
\lecturetimestamp{02:02:21--02:03:17}}
\end{classroomqa}

Теперь две половины рассуждения сходятся. Измерение имеет смысл только после того, как заданы его управляющие параметры и условия равновесия. Приближение имеет смысл только после того, как заданы его малый параметр и критерий неприменимости. Термодинамическая величина является функцией состояния лишь после того, как её дифференциал проходит проверку на независимость от пути. Во всех случаях строгая математика не заменяет физическую интерпретацию; она точно указывает, где эта интерпретация оправдана.
